\section{INTRODUCTION}

Selenium is an essential trace element for mammals that participates in the biosynthesis of selenoproteins such as glutathione peroxidase, playing a pivotal role in maintaining organismal redox homeostasis. Compared to inorganic selenium species such as selenite and selenate, zero-valent selenium nanoparticles (SeNPs) exhibit higher bioavailability and lower acute toxicity. Benefiting from their favorable biocompatibility and diverse biological activities, SeNPs hold broad application prospects across biomedical fields, including anti-tumor, antibacterial, antidiabetic, and antioxidant therapeutics \cite{hosnedlova2018nanoselenium,bisht2022senps,sowmya2024nanoselenium,karthik2024senps}. However, selenium possesses a narrow therapeutic index with a limited margin between effective and toxic doses; even within pharmacologically active dose regimens, SeNPs can still induce adverse outcomes such as oxidative damage and cytotoxicity \cite{wang2020prooxidation}. Furthermore, excessive exposure may precipitate systemic toxicities, including fever, nausea, respiratory, cardiovascular, and gastrointestinal impairments, and even renal failure \cite{hosnedlova2018nanoselenium}, highlighting the necessity and urgency of comprehensive biosafety evaluations.

The biotoxicity of nanomaterials is collectively governed by multifaceted factors such as particle size, morphology, surface properties, and chemical composition. Marked non-linear coupling characteristics exist among these physicochemical parameters, compounding the complexity of toxicity prediction and mechanistic elucidation \cite{ji2025rosml}. Conventional nanotoxicity assessments predominantly rely on in vitro cell-based assays and in vivo animal models (e.g., rat, mouse, and zebrafish models); nevertheless, such experimental paradigms generally suffer from protracted testing cycles, prohibitive costs, and constrained throughput \cite{ahmadi2024toxicity}. With the advancement of computational toxicology, nano-quantitative structure--activity/toxicity relationship (nano-QSAR/QSTR) methodologies have progressively emerged as indispensable complementary tools in nanotoxicity research \cite{roy2023tio2}. In recent years, machine learning algorithms such as Random Forest (RF), Support Vector Machines (SVM), and Neural Networks (NN) have been widely deployed to predict the cytotoxicity of metal and metal oxide nanoparticles \cite{zhang2023qsar}, frequently coupled with clustering analysis and feature selection techniques to further enhance model stability, predictive accuracy, and generalizability \cite{sang2022mlcytotoxicity}.

Dedicated investigations on selenium nanoparticles have revealed that their cytotoxic effects are intimately linked to factors including particle size, exposure concentration, and surface functionalization. Regarding particle size, as SeNP dimensions decrease and specific surface area increases, cellular uptake efficiency and biological activity are typically amplified, accompanied by an elevated risk of potential toxicity \cite{zhang2023senpsbio}. In terms of concentration, Wang et al. \cite{wang2020prooxidation} demonstrated that exposing rat hepatocytes to 24~$\mu$M and 48~$\mu$M SeNPs for 24~h reduced cell viability to 87.76\% and 51.37\%, respectively, manifesting an evident dose-dependent cytotoxic response. Surface modification likewise exerts a substantial impact on the biosafety and toxicological behavior of SeNPs. For instance, surface functionalization with polysaccharides such as Spirulina polysaccharide (SPS), Gracilaria lemaneiformis polysaccharide (GLP), and Astragalus polysaccharide (AASP) can potentiate the selective killing efficacy of SeNPs against tumor cells while mitigating their toxicity toward normal cells \cite{zhang2023senpsbio,ji2023astragalus}; conversely, quaternized chitosan--gallic acid-modified SeNPs exhibit pronounced cytotoxicity in Colo357 and HaCaT cells \cite{lunkov2023chitosan}. These findings substantiate that the toxicological profile of SeNPs exhibits pronounced physicochemical dependency that defies simplistic empirical extrapolation, underscoring the necessity of establishing toxicity prediction models tailored to specific physicochemical descriptors of selenium nanoparticles.

Recently, machine learning and deep learning have exhibited salient strengths in high-dimensional feature mining, complex non-linear relationship modeling, and pattern recognition, offering innovative technological routes for computational toxicology \cite{hosni2025mlnp}. Shirokii et al. utilized correlation-based feature selection and grid search hyperparameter optimization to elevate the 10-fold cross-validation $Q^2$ of a quantitative cytotoxicity prediction model for inorganic nanomaterials to 0.86 \cite{shirokii2023quantitative}. Furthermore, ensemble learning architectures such as LightGBM have displayed superior predictive performance over conventional machine learning models in distinguishing cytotoxic from non-cytotoxic nanoparticles \cite{hosni2025mlnp}, while hybrid modeling frameworks such as Boruta-RF-ARM have further pinpointed critical toxicological determinants, including exposure concentration, core--shell composition, and cell lineage \cite{yildirim2025rheumatoid}. Despite these advancements in nanotoxicity modeling, prevailing studies have largely focused on organic small molecules or generalized nanomaterial categories. Dedicated investigations into the cytotoxic profile of selenium nanoparticles remain scarce, and domain-specific predictive models alongside feature engineering frameworks for SeNPs are still lacking. To bridge this gap, this study establishes an analytical framework for selenium nanoparticle cytotoxicity prediction comprising multi-model benchmarking, enhanced feature engineering, and automated feature selection. Taking essential variables—such as concentration, particle size, surface modification, and material chemical properties—as inputs, this work systematically evaluates the applicability and predictive efficacy of diverse models across both generalized nanoparticle benchmarks and dedicated selenium nanoparticle datasets, aiming to provide a methodological foundation for SeNP biosafety assessment, experimental candidate prioritization, and the future development of intelligent prediction platforms.

% 硒是哺乳动物必需的微量元素，参与谷胱甘肽过氧化物酶等硒蛋白的合成，在维持机体氧化还原稳态中发挥重要作用。与亚硒酸盐、硒酸盐等无机硒相比，零价态纳米硒具有更高的生物利用度和更低的急性毒性。得益于其良好的生物相容性及多种生物学活性，纳米硒在抗肿瘤、抗菌、抗糖尿病及抗氧化等生物医学领域展具有广阔的应用前景\cite{hosnedlova2018nanoselenium,bisht2022senps,sowmya2024nanoselenium,karthik2024senps}。然而，硒的治疗指数（therapeutic index）较低，其有效剂量与毒性剂量之间界限有限，即使在药理活性剂量范围内，纳米硒仍可能产生氧化损伤及细胞毒性等不良反应\cite{wang2020prooxidation}。过量暴露还可能导致发热、恶心，以及呼吸系统、心血管系统、消化系统损伤甚至肾衰竭等系统性毒性反应\cite{hosnedlova2018nanoselenium}，显示其生物安全性评价的重要性与必要性。
 
% 纳米材料的生物毒性受到粒径、形貌、表面性质及化学组成等多种因素的综合影响，各理化参数之间存在显著的非线性耦合特征，增加了毒性效应预测与机制解析的复杂性\cite{ji2025rosml}。传统的纳米毒性评价主要依赖体外细胞实验以及体内动物实验（如大鼠、小鼠和斑马鱼模型），但该类方法普遍存在实验周期长、成本高及检测通量有限等问题\cite{ahmadi2024toxicity}。随着计算毒理学的发展，纳米定量构效关系（nano-QSAR/QSTR）等方法逐渐成为纳米毒性研究的重要补充工具\cite{roy2023tio2}。近年来，随机森林（RF）、支持向量机（SVM）及神经网络（NN）等机器学习算法已被广泛应用于金属及金属氧化物纳米颗粒的细胞毒性预测\cite{zhang2023qsar}，并结合聚类分析、特征筛选等方法进一步提升模型的稳定性、预测精度及泛化能力\cite{sang2022mlcytotoxicity}。
 
% 针对纳米硒的研究表明，其毒性效应与粒径、暴露浓度及表面修饰等因素密切相关。粒径方面，随着纳米硒颗粒粒径减小及比表面积增加，其生物活性和细胞摄取效率通常进一步增强，同时潜在毒性风险也随之升高\cite{zhang2023senpsbio}。浓度方面，Wang 等\cite{wang2020prooxidation}的研究表明，大鼠肝细胞暴露于 24~$\mu$M 和 48~$\mu$M SeNPs 24~h 后，细胞存活率分别下降至 87.76\% 和 51.37\%，呈现明显的剂量依赖性毒性效应。表面修饰同样能够显影响纳米硒的生物安全性与毒性行为。例如，螺旋藻多糖（SPS）、江蓠多糖（GLP）及黄芪多糖（AASP）等表面修饰可增强纳米硒对肿瘤细胞的选择性杀伤作用，并降低其对正常细胞的毒性\cite{zhang2023senpsbio,ji2023astragalus}；而季铵化壳聚糖--没食子酸修饰的 SeNPs 在 Colo357 和 HaCaT 细胞中则表现出较高的细胞毒性\cite{lunkov2023chitosan}。上述研究表明，纳米硒的毒性效应具有显著的理化性质依赖特征，难以通过传统经验进行简单外推，因此有必要建立面向纳米硒特定物化参数的毒性预测模型。
 
% 近年来，机器学习与深度学习在高维特征挖掘、复杂非线性关系建模及模式识别等方面展现出显著优势，为计算毒理学研究提供了新的技术路径\cite{hosni2025mlnp}。Shirokii 等通过相关性特征筛选及网格搜索优化超参数，将无机纳米材料细胞毒性的定量预测模型在 10 折交叉验证中的 $Q^2$ 提升至 0.86\cite{shirokii2023quantitative}。此外，LightGBM 等集成学习算法在区分细胞毒性与非细胞毒性纳米颗粒方面表现出优于传统机器学习模型的预测性能\cite{hosni2025mlnp}
，Boruta-RF-ARM 等组合建模框架则进一步识别出浓度、核壳材料及细胞类型等关键毒性影响因素\cite{yildirim2025rheumatoid}。尽管现有研究已在纳米毒性预测领域取得一定进展，但相关模型主要集中于有机小分子或通用纳米材料，对纳米硒毒性效应的针对性研究仍相对有限，其专用预测模型与特征建模方法尚不完善。基于此，本研究围绕纳米硒毒性预测构建了由多模型比较、增强特征建模及自动特征筛选组成的分析框架，以浓度、粒径、表面修饰及材料化学属性等关键变量作为输入，系统评估不同模型在通用纳米颗粒数据集与纳米硒专用数据集中的适用性与预测性能，旨在为纳米硒生物安全性评价、实验优先级筛选及后续智能预测平台开发提供方法学基础。
